Your task is to calculate the area of a given polygon.
The polygon consists of $n$ vertices $(x_1,y_1),(x_2,y_2),\dots,(x_n,y_n)$. The vertices $(x_i,y_i)$ and $(x_{i+1},y_{i+1})$ are adjacent for $i=1,2,\dots,n-1$, and the vertices $(x_1,y_1)$ and $(x_n,y_n)$ are also adjacent.

\vspace{0.3cm}\noindent\textbf{Input}

The first input line has an integer $n$: the number of vertices.
After this, there are $n$ lines that describe the vertices. The $i$th such line has two integers $x_i$ and $y_i$.
You may assume that the polygon is simple, i.e., it does not intersect itself.

\vspace{0.3cm}\noindent\textbf{Output}

Print one integer: $2a$ where the area of the polygon is $a$ (this ensures that the result is an integer).

\vspace{0.3cm}\noindent\textbf{Constraints}

\textbullet\ $3 \le n \le 1000$
\textbullet\ $-10^9 \le x_i, y_i \le 10^9$